% Generated draft: requires mathematical and source review.
\documentclass[11pt]{article}
\usepackage[margin=25mm]{geometry}
\usepackage{fontspec}
\setmainfont{Latin Modern Roman}
\usepackage{amsmath,amssymb,mathtools,unicode-math}
\setmathfont{Latin Modern Math}
\usepackage{enumitem,needspace,xurl,hyperref}
\setlength{\parindent}{0pt}
\setlength{\parskip}{5pt}
\setlength{\emergencystretch}{3em}
\newcommand{\N}{\mathbb{N}}
\newcommand{\Z}{\mathbb{Z}}
\newcommand{\Q}{\mathbb{Q}}
\newcommand{\R}{\mathbb{R}}
\newcommand{\C}{\mathbb{C}}
\newcommand{\F}{\mathbb{F}}
\newcommand{\A}{\mathbb{A}}
\newcommand{\PP}{\mathbb{P}}
\newcommand{\E}{\mathbb{E}}
\newcommand{\eps}{\varepsilon}
\newcommand{\dd}{\,\mathrm{d}}
\begin{document}
\section*{UnsolvedMath ID: 30005802}
\textbf{Source problem code:} catalogue entry
\subsection{Definitions and mathematical statement}
This record contains three separate questions. The definitions below make conventional interpretations explicit, but the supplied text does not determine all hypotheses needed to identify the intended conjectures uniquely.

\begin{enumerate}
\item A finite geometric simplicial complex in a Euclidean space is a finite collection of simplices closed under taking faces, such that the intersection of any two simplices is a face of each. Its support \(|K|\) is the union of its simplices. Under this convention, a triangulation of a set \(X\) is such a complex with support \(X\).

A stellar subdivision at a nonempty face \(\sigma\) inserts a point \(v\) in the relative interior of \(\sigma\) and replaces the closed star \(\sigma*\operatorname{lk}_K(\sigma)\) by
\[
v*\partial\sigma*\operatorname{lk}_K(\sigma),
\]
leaving the rest of the complex unchanged. Here \(*\) denotes simplicial join, \(\partial\sigma\) is the boundary complex, and \(\operatorname{lk}_K(\sigma)\) consists of those faces \(\eta\) disjoint from \(\sigma\) for which \(\eta\cup\sigma\) is a face of \(K\). An iterated stellar subdivision is obtained by a finite sequence of such operations, allowing a sequence of length zero.

The first supplied question, labeled ``Oda's factorization conjecture,'' asks whether
\[
|K|=|L|\quad\Longrightarrow\quad
\exists M\;\bigl(K\rightsquigarrow M\ \text{and}\ L\rightsquigarrow M\bigr),
\]
where each arrow denotes an iterated stellar subdivision and \(M\) is the same resulting triangulation. The intended class of triangulations, the ambient dimension, and any rationality, lattice, fan, smoothness, or restrictions on permitted stellar subdivisions are not specified. Thus this formula records the common-subdivision question within an unspecified admissible class; the unrestricted geometric convention above must not be treated as a verified identification of the named conjecture.

\item For a finite set \(A\subseteq\R^d\), one conventional meaning of a triangulation of \(A\) is a geometric simplicial complex with support \(\operatorname{conv}(A)\) whose vertices belong to \(A\). Whether every point of \(A\) must occur as a vertex is an additional convention. A circuit is an inclusion-minimal affinely dependent subset of \(A\). In the usual geometric setting, a bistellar flip replaces a local triangulation associated with one of the two triangulations of a circuit by the other, with a common link, provided the replacement gives another permitted triangulation.

The supplied question literally asks whether ``any two point sets'' in \(\R^3\) and \(\R^4\) are connected by bistellar flips. It does not specify states or admissible moves for a graph whose vertices are point sets. A possible intended question is the following fixed-configuration statement:
\[
\forall d\in\{3,4\}\ \forall A\subseteq\R^d\ \text{finite}\quad
\forall T,T'\in\mathcal T(A)\quad
\exists m\geq0\ \exists T_0,\ldots,T_m\in\mathcal T(A)
\]
\[
T_0=T,\qquad T_m=T',\qquad
T_i\ \text{and}\ T_{i+1}\ \text{differ by one permitted bistellar flip}
\quad(0\leq i<m).
\]
Here \(\mathcal T(A)\) would have to be the intended class of triangulations of a fixed configuration. This is a conditional reconstruction, not a correction established by the supplied evidence. Finiteness, full dimensionality, general position, use of every point, and the precise flip convention remain unspecified.

\item A lattice polytope in \(\R^d\) is the convex hull of finitely many points of \(\Z^d\). For a positive scalar \(c\), write \(cP=\{cx:x\in P\}\). For a lattice polytope \(Q\), put \(k=\dim Q\) and
\[
\Lambda_Q=\operatorname{span}_{\R}(Q-Q)\cap\Z^d.
\]
A lattice \(k\)-simplex \(\operatorname{conv}(v_0,\ldots,v_k)\) in the affine span of \(Q\) is unimodular relative to this induced lattice if \(v_1-v_0,\ldots,v_k-v_0\) form a \(\Z\)-basis of \(\Lambda_Q\). A unimodular triangulation of \(Q\) is a triangulation with lattice vertices whose maximal simplices have this property. In the full-dimensional case, the condition is
\[
\left|\det(v_1-v_0,\ldots,v_d-v_0)\right|=1.
\]

The third question asks whether, for every dimension \(d\), there is a positive dilation factor \(c_d\), depending only on \(d\), such that every lattice polytope \(P\subseteq\R^d\) has a unimodular triangulation after dilation:
\[
\forall d\geq1\ \exists c_d>0\ \forall P\subseteq\R^d\ \text{a lattice polytope},
\qquad c_dP\ \text{admits a unimodular triangulation}.
\]
The positivity of the factor is the nondegenerate dilation interpretation. The record does not explicitly state whether \(c_d\) must be an integer or whether \(P\) must be full dimensional. The relative-lattice convention above explains one way to include lower-dimensional polytopes; its intended use requires confirmation. No regularity requirement for the triangulation is stated.
\end{enumerate}
\subsection{Short English statement}
The record asks three questions: whether two triangulations with the same support have a common iterated stellar subdivision; whether a suitably specified bistellar-flip graph is connected in dimensions \(3\) and \(4\); and whether a dilation factor depending only on dimension guarantees unimodular triangulations of all lattice polytopes. The admissible objects in the first question and the meaning of the second question require clarification. The third also leaves some dilation and dimension conventions unstated. The supplied open-status assessment is dated curation, not independent verification.
\subsection{Sources}
\begin{itemize}
\item Local UnsolvedMath registry, supplied complete record with permanent ID \(30005802\), title ``Three Problems on Triangulations,'' source problem code \texttt{catalogue entry}, Oberwolfach Reports Open Problems collection. The supplied proposed year is 2024; no proposer is identified.
\item Supplied report citation: \emph{Discrete Geometry} (2024), Oberwolfach Reports, pp.~137--202, DOI \texttt{10.4171/owr/2024/3}. Source URL and supplied literature-source link, labeled ``Discrete Geometry (2024)'': \url{https://doi.org/10.4171/owr/2024/3}.
\item The source labels the first question ``Oda's factorization conjecture.'' The supplied excerpt does not identify the speaker or contributor responsible for the three-question list. The subsequent heading ``A Complex Analogue of the Goodman-Pollack-Wenger Theorem'' and name Daniel McGinnis occur after that list and do not establish attribution of these questions.
\item Literature assessment date: \texttt{2026-08-22}. The supplied dated curation labels the record open, describes it as bundling three distinct open problems, and states that the report presents them as open and that no single theorem resolves all three. This is supplied dated curation, not independently verified literature assessment. Its instruction to verify time-sensitive status and cited claims independently is retained.
\item The registry supplies statement status \texttt{corrected\_verified} and states that the statement was reconstructed and checked against the official Oberwolfach report. These are supplied provenance claims; this draft does not independently verify them or the cited report.
\end{itemize}
\end{document}
